\chapter{Literature Review}
\label{chap:literature}

% YOU WRITE THIS CHAPTER. Ask me to explain any paper/concept below before
% you write about it -- I can summarise a paper's actual contribution for
% you to then describe in your own words, but I won't write the review text.

\section{Prompting Strategies for Reasoning}
\subsection{Chain-of-Thought Prompting}
% \cite{wei2022chain} -- ask me what this paper actually showed and why
% it matters as your baseline condition.

\subsection{Tree-of-Thought Prompting}
% \cite{yao2023tree} -- ask me to explain the full tree-search design vs.
% the reduced-scale multi-branch version actually implemented here, and
% why that's a defensible scope decision to state explicitly, not hide.

\section{Mathematical Reasoning Benchmarks}
\subsection{The MATH Dataset}
% \cite{hendrycks2021measuring}

\subsection{AIME as a Benchmark Source}
% \cite{maa_aime} -- also: why AIME for the Hard tier specifically, and
% the data-contamination risk of any public competition-math benchmark.

\section{LLM-as-Judge Evaluation}
% \cite{zheng2023judging} -- and the self-evaluation bias risk this
% specific project has (Gemini judges Gemini's own responses), and how
% the neutral-judge validation sample addresses it.

\section{Error Analysis of LLM Mathematical Reasoning}
% Needs an actual literature search -- ask me to help you find and
% summarise papers that classify or localise LLM reasoning errors
% specifically (not just accuracy benchmarks), so you can position this
% thesis's contribution against them.

\section{Summary}
% Tie the sections above together into "here's the gap, here's what
% this thesis does about it."
